	\section{Priority Assignments Optimization}
    \label{section_priority_opt}

    This section first addresses a simplified version of the SP optimization problem, the \emph{priority-assignment subproblem} with the QoS budgets $\boldsymbol{Q}$ held \emph{fixed}:
    \begin{equation}
        \max_{\boldsymbol{P}} \textbf{SP}(\boldsymbol{P};\boldsymbol{Q}, \textbf{E})
        \label{eq_pa_subproblem}
    \end{equation}
	\begin{equation}
		Pr(r_i > D_i) \leq \Theta_i, \quad \forall \tau_i \in \boldsymbol{\tau}^{safe}
		\label{eq_important_task_constraint2}
	\end{equation}
    Section~\ref{section_config_opt} lifts the QoS-fixed restriction and optimizes the priority assignment $\boldsymbol{P}$ and the QoS budgets $\boldsymbol{Q}$ jointly.

    Priority assignment is a fundamental scheduling technique whose effectiveness directly governs the safety and performance of real-time robotic systems such as autonomous vehicles and aerospace platforms.
    The remainder of this section develops an efficient online optimizer: we start from a brute-force baseline, refine it with a Constrained Audsley Beam Search (CABS) built on Audsley's algorithm~\cite{Audsley91optimalpriority}, and then introduce incremental optimization techniques that trade a bounded amount of optimality for the run-time efficiency required online.


	\subsection{Brute-Force Priority Assignments}
	\label{section_bf_pa}
    Given a task set of $n$ tasks, the brute-force algorithm enumerates and evaluates all $n!$ possible priority assignment vectors to identify the optimal solution.
    Each enumerated assignment is checked against \eqref{eq_important_task_constraint2}, and any that violates an important task's threshold is discarded before scoring, 
    so the optimum is taken over feasible assignments only. 
    It guarantees to find the global optimal solution if execution time prediction is accurate.
    While straightforward, this approach has an exponential computational complexity and is not practical in reality.
    We use it as an offline optimality reference in our experiments (Section~\ref{section: simulation}).

	\subsection{Constrained Audsley Beam Search}
	\label{section_cabs}
    Since the QoS budgets are fixed, the system-level SP metric is highly correlated with the system's schedulability.
    This motivates us to build on the classical Audsley's algorithm~\cite{Audsley91optimalpriority}, which is optimal in schedulability under certain conditions~\cite{Audsley91optimalpriority} and has $O(n^2)$ runtime complexity, where $n$ is the number of tasks.
    Directly applying Audsley's algorithm has two limitations: it optimizes schedulability rather than the SP metric, so its assignments need not maximize SP, and it assumes hard real-time constraints.

    To address these, we propose the \emph{Constrained Audsley Beam Search}, which extends Audsley's iterative structure along two axes:
    \begin{itemize}
        \item \textbf{Beam.} Retain multiple ($b$) partial priority assignments per iteration rather than a single greedy choice, scored by accumulated SP-value loss rather than schedulability alone.
        \item \textbf{Constraint.} Reject any partial assignment that violates a safety-critical task's deadline-miss threshold, generalizing Audsley's hard-real-time feasibility test to the probabilistic bound $Pr(r_i > D_i) \leq \Theta_i$.
    \end{itemize}
    CABS preserves Audsley's $O(n^2)$ iterative structure.
    The algorithm works iteratively from the first iteration to the \ith{N} iteration where \ith{N} is the number of tasks in the task set $\boldsymbol{\tau}$.
    In the \ith{k} iteration, we try to find the task $\tau_i$ that can be assigned with the \ith{k} lowest priority based on the following rules:
    \begin{itemize}
        \item $\tau_i$ is not assigned a priority value yet. Each task will only be assigned priority once and will not be changed while other tasks are assigned priorities;
        \item If $\tau_i$ is a safety-critical task, assigning $\tau_i$ the \ith{k} lowest priority will not violate its deadline-miss threshold;
        \item Among all the tasks $\boldsymbol{\tau}^{U}$ whose priorities are not decided yet, $\tau_i$ is the task with minimum regret (SP value loss):
    \end{itemize}
    \begin{equation}
    \begin{split}
        \tau_i = \argmin_{\tau_i \in \boldsymbol{\tau}^{U}} \big( 1.0 & - \mathcal{P}(\tau_i; \textbf{E}) \\
        & \cdot \mathcal{S}\!\left(Pr(r_i > D_i ; \boldsymbol{P}(\boldsymbol{\tau}^{U}), \boldsymbol{Q}, \textbf{E}) - \Theta_i\right) \big)
    \end{split}
    \label{eq:cabs_goal}
    \end{equation}
    where $\boldsymbol{P}(\boldsymbol{\tau}^{U})$ denotes the priority assignments of tasks in $\boldsymbol{\tau}^{U}$.
    
This heuristic allows the algorithm to balance schedulability with SP metric optimization.
The pseudocode is shown below:

	\begin{algorithm}[htbp]
		\caption{Constrained Audsley Beam Search (CABS)} \label{alg:cabs}
		\KwIn{Task set $\boldsymbol{\tau}$, beam width $b$}
		\KwOut{Selected priority assignment}
		$\partialpa = \{$ a single empty partial path $p$ with $p.\textsc{Pool} \leftarrow \boldsymbol{\tau}$ $\}$ \tcp{lowest-priority task appended first}
		\For{$k$ in $0,...,N-1$}
		{
			$\partialpacur=\{\emptyset\}$\\
			\For{$p$ in $\partialpa$}{
				\For{$\tau_i$ in $p.\textsc{Pool}$}{
					$p' \leftarrow \textbf{copy}(p)$ \tcp{copy before extending}
					\If{$p'.\textsc{Assign}(\tau_i)$}{
					$\partialpacur.\textbf{Push}(p')$ \tcp{append $\tau_i$ at priority $k$, accumulate SP loss; reject (return false) if $\tau_i \in \boldsymbol{\tau}^{safe}$ and $Pr(r_i > D_i) > \Theta_i$}
					}
				}
			}
			$\partialpa = \textbf{SelectTop}(\partialpacur, b)$ \tcp{keep the $b$ paths of \emph{lowest} accumulated SP loss}
		}
		\KwRet $\textbf{SelectTop}(\partialpa, 1)$
	\end{algorithm}

	\begin{Example}
        Consider a task set of three tasks: $\tau_0$, $\tau_1$, and $\tau_2$.
        Assume we retain two partial priority assignment vectors during each iteration.

        First iteration.
        Assign the lowest priority to each task and evaluate the SP loss.
        Suppose $\tau_0$ and $\tau_1$ have the smallest SP loss.
        The retained partial priority assignments are:
        \begin{itemize}
        \item $P_0$: $\{\tau_0 \}$.
        \item $P_1$: $\{\tau_1 \}$.
        \end{itemize}

        Second iteration.
        For each retained partial priority assignment, assign the second-lowest priority to the remaining tasks and evaluate the accumulated SP loss:
        \begin{itemize}
            \item $P_2: \{\tau_0, \tau_1\} = 0.5$
            \item $P_3: \{\tau_0, \tau_2\} = 0.4$
            \item $P_4: \{\tau_1, \tau_0\} = 0.6$
            \item $P_5: \{\tau_1, \tau_2\} = 0.8$
        \end{itemize}
        Therefore, the two partial priority assignments to keep after the second iteration are $P_2$ and $P_3$.

        Third iteration.
        Assign the third-lowest priority to the remaining task for each retained assignment:
        \begin{itemize}
            \item $P_6:  \{\tau_0, \tau_1, \tau_2\}$
            \item $P_7:  \{\tau_0, \tau_2, \tau_1\}$
        \end{itemize}
        Evaluate the SP loss for these two assignments and return the one with the minimum SP loss as the final solution.
	\end{Example}


	\subsection{Incremental Priority Assignment Optimization}
	\label{section_increment_pa}
    Since the environment changes continuously and the scheduler is triggered frequently with short time intervals, it's reasonable to expect tasks' configurations and their optimal priority assignments do not change dramatically. 
    Therefore, we can leverage the last solution found to perform efficient incremental optimization by focusing only on the tasks whose execution time distribution change.
    
    The incremental optimization always starts from an initial feasible solution (such as one found by the Constrained Audsley Beam Search).
    Starting from this initial solution, assume there are $n$ tasks whose execution time become different since the last time the scheduler is called, and we sort them based on their importance (SP weights). 
    The incremental optimization iterates through these $n$ tasks one by one: for each task $\tau_i$, we try to adjust its priority value while keeping all other tasks' relative priority order unchanged. 
    Assume the task set has $N$ tasks in total, there will be $N$ options for $\tau_i$'s new priority level. So we follow a simple heuristic based on improving task set schedulability:
The possible scenarios are:
\begin{enumerate}
    \item If $\tau_i$ execution time increases, the re-search direction is \emph{downward} for improving overall schedulability.
    \item If $\tau_i$ execution time decreases, the re-search direction is \emph{upward} for explorin improving performance.
    \item If $\tau_i$ has the uniquely highest weight, we do not want to starve this task, the re-search direction is \emph{upward}.
\end{enumerate}
In every scenario $\tau_i$ is removed from the current priority vector and re-inserted at each position within the chosen half; 
among the positions that preserve \eqref{eq_important_task_constraint2}, the one yielding the highest SP metric is adopted if it strictly improves on the current assignment, 
otherwise the current position is retained.
Positions that would drive an important task over its threshold are skipped, so the incremental search, like CABS, never moves to an infeasible assignment.

Potential issues in terms of optimal solution drift caused by the incremental optimization is mainly solved by periodically triggering CABS, see Section~\ref{sec: drift}.


